% ID: SEL-D0034-OA-E04
% ESTADO: validada. S035, cotejada con examen y criterios.
\textbf{a)} La proporción observada es $\hat p=19/2000=0.0095$.
Con 19 defectos y 1981 artículos sin defecto, usamos la
aproximación normal bilateral:
\[
0.0095\pm1.959964\sqrt{\frac{0.0095(0.9905)}{2000}}
\approx\boxed{[0.00525,\ 0.01375]}.
\]
El intervalo va aproximadamente del $0.525\%$ al $1.375\%$.
El $95\%$ se refiere a la cobertura del procedimiento:
al repetir el muestreo, aproximadamente el $95\%$ de los
intervalos así construidos contendrían la proporción real.

\textbf{b)} Interpretamos $1\%$ como un punto porcentual y
usamos $0.0095$ como estimación previa de la proporción.
En el mismo modelo normal, al $99\%$ se requiere
\[
n\ge z_{0.995}^2\frac{0.0095(0.9905)}{0.01^2}
\approx624.3272.
\]
El mínimo planificado es $\boxed{625}$ artículos. Con esa
proporción se esperan unos seis defectos, por lo que el margen
es aproximado. Si se planifica para el peor caso sin estimación
previa, $p(1-p)\le1/4$ da 16588 artículos.
